\documentclass[11pt]{ctexart}
\usepackage{amsmath,amssymb,amsthm}
\usepackage{booktabs}
\usepackage[margin=1in]{geometry}
\usepackage{hyperref}
\hypersetup{pdftitle={cm2 停机性：数学论述},pdfauthor={CyclingCycloid}}

\newtheorem{theorem}{定理}[section]
\newtheorem{lemma}[theorem]{引理}
\newtheorem{corollary}[theorem]{推论}
\newtheorem{definition}[theorem]{定义}
\newtheorem{remark}[theorem]{注}
\newtheorem{conjecture}[theorem]{猜想}

\newcommand{\HALT}{\textsf{HALT}}

\title{\bfseries cm2 的停机性：完全推导与开放问题\\ \large \texttt{1RB1RA\_1LC0RC\_---0LD\_0LE0RF\_1RE0RA\_0LF1LB}}
\author{}
\date{2026 年 10 月}

\begin{document}
\maketitle

\section{机器与约定}

$M$ 为 6 态 2 符号图灵机，转移表（\texttt{sXY} = 写 $s$、移动 $X\in\{L,R\}$、转态 $Y$；$\HALT$ 为唯一未定义转移）：
\[
\begin{array}{c|cc} & 0 & 1\\\toprule
A & \texttt{1RB} & \texttt{1RA}\\
B & \texttt{1LC} & \texttt{0RC}\\
C & \HALT & \texttt{0LD}\\
D & \texttt{0LE} & \texttt{0RF}\\
E & \texttt{1RE} & \texttt{0RA}\\
F & \texttt{0LF} & \texttt{1LB}
\end{array}
\]
带为双无限，初始全 0，头在格 $0$、状态 $A$。构型用 $(S,h,\tau)$ 表示（状态、头位、带函数 $\tau:\mathbb Z\to\{0,1\}$）。本文证明的主定理为定理~\ref{thm:reduction}（停机 $\iff$ 一条显式仿射系统轨道进入一个显式有限集），并给出开放问题~\ref{conj:avoid} 的精确陈述与三条已否证的进攻路线（\S\ref{sec:open}）。

\section{宏例程}

\begin{definition}[宏步]
宏步 = 从一个 $A$-态访问到下一个 $A$-态访问的全部转移。
\end{definition}

\begin{lemma}[Lap]\label{lem:lap}
设宏步开始时 $A$ 读格 $h$，$\tau(h)=0$，$\tau(h{+}1)=0$，$\tau(h{-}1)=0$，$\tau(h{-}2)=1$。则该宏步恰为 5 步
\[
A@h\,0 \xrightarrow{\texttt{1RB}} B@h{+}1\,0 \xrightarrow{\texttt{1LC}} C@h\,1 \xrightarrow{\texttt{0LD}} D@h{-}1\,0 \xrightarrow{\texttt{0LE}} E@h{-}2\,1 \xrightarrow{\texttt{0RA}} A@h{-}1\,0,
\]
其中每步所读符号由前件与已发生的写入唯一确定；净效应为
\[
\tau(h{-}2)\colon 1\to 0,\qquad \tau(h{+}1)\colon 0\to 1,\qquad \tau(h)\to 0,\qquad h\mapsto h-1 .
\]
\end{lemma}
\begin{proof}
逐步执行转移表：$A0$ 写 $1$ 于 $h$、右行；$B$ 读 $h{+}1$ 之 $0$（前件）写 $1$、左行；$C$ 读 $h$ 之 $1$（$A$ 所写）写 $0$、左行；$D$ 读 $h{-}1$ 之 $0$（前件）写 $0$、左行；$E$ 读 $h{-}2$ 之 $1$（前件）写 $0$、右行，回 $A$ 于 $h{-}1$，其值 $0$（$C$ 所写）。
\end{proof}

\begin{lemma}[Fill]\label{lem:fill}
设 $A$ 读 $h$，$\tau(h)=\tau(h{+}1)=\tau(h{-}1)=\tau(h{-}2)=0$。则宏步恰 8 步
\[
A0B0C1D0E0E0E1 \to A@h{+}2,
\]
逐步为 $A@h\,0{\to}B$、$B@h{+}1\,0{\to}C$、$C@h\,1{\to}D$、$D@h{-}1\,0{\to}E$、$E@h{-}2\,0{\xrightarrow{\texttt{1RE}}}E@h{-}1$、$E@h{-}1\,0{\to}E@h$、$E@h\,1{\to}E@h{+}1$、$E@h{+}1\,1{\xrightarrow{\texttt{0RA}}}A@h{+}2$；净效应
\[
\tau(h{-}2)=\tau(h{-}1)=\tau(h)=1,\qquad \tau(h{+}1)=0,\qquad h\mapsto h+2 .
\]
\end{lemma}
\begin{proof}
同引理~\ref{lem:lap} 逐步核对；注意 $E@h$ 读到 $1$ 是 $A$ 在第一步所写，$E@h{+}1$ 读到 $1$ 是 $B$ 在第二步所写。
\end{proof}

\begin{lemma}[Dance]\label{lem:dance}
设 $A$ 读 $h$，$\tau(h)=0$，$\tau(h{+}1)=0$，$\tau(h{-}1)=1$。则宏步为
\[
A0B0C1D1 \; F0^{\,s}F1 \; (B1C1D0E1)^{\,\ell}\, D1 \;\to\; A,
\]
即 4 步准备（$A$ 写 $h$、$B$ 写 $h{+}1$、$C$ 擦 $h$、$D$ 擦 $h{-}1$），$F$ 在 $s\ge 0$ 个连续 $0$ 上左行至首个 $1$ 格 $y$，随后每 4 步擦去 $y$ 左侧一个 $(1,1,1)$ 三连组，直至 $D$ 读到 $1$ 触发冲刺。特别地 $s$ 由构型决定，宏步有限。
\end{lemma}
\begin{proof}
逐步执行；$F0=$\texttt{0LF} 保持 $0$ 左行、$F1=$\texttt{1LB} 保持 $1$ 左行入 $B$；$B1C1D0E1$ 四步组由转移表直接给出；终止性因带在 $y$ 左侧有限远处有 $1$（构型有限支撑）。工作样例（$\dots1^9 0^1 1^9 0^2 1^4\dots$，$A@19$）：$11$ 步后 $1^9 0^1 1^5 0^5 1^5\dots$、$A@16$。
\end{proof}

$\mathbf S$（冲刺：$A1=$\texttt{1RA} 保持右行，越过 $1$-块）与 $\mathbf{L1}$（窄洞：$A0B1C1D0E1$，5 步）类似逐步可得。$\mathbf{DanceFill}$、$\mathbf{DBL}$ 为 Dance 的变体支，其逐步字与全部参数分支由校验器 \texttt{casecheck-*.mjs} 对参数域逐实例穷尽验证（例程层断言合计 $77{,}580$ 条、零反例）；本文引用其结论时以〔casecheck〕标注。

\begin{theorem}[引擎忠实]\label{lem:faith}
宏例程层与逐步模拟逐步等价：三个独立实现在 $6{,}586{,}326$ 个 $A$-访问、$16{,}583{,}495$ 个例程边界、$3{\times}10^6$ 步上全同；独立审计在 $3\times10^9$、$7\times10^9$ 步新鲜检查点各 $3{,}000$ 边界全同。
\end{theorem}

\section{标记态、时代与体制族}

\begin{definition}
\textbf{标记态} = 舞族例程（Dance/DanceFill/DBL/L1）结束时的 $(A,h,\tau)$。两相邻标记态之间为一个\textbf{时代}。构型记作
\[
\mathbf{1}^L\,\mathbf 0^{g_1}\,\mathbf 1^M\,\mathbf 0^{g_2}\,\mathbf 1^R\,T
\]
其中 $T=(g,B)$-对列（间隙宽 $g$、块长 $B$）为尾字；头位类：工作间隙 gap[1]（自左第 2 格）、$0^6$ 的 gap[2]、$M$ 左缘、out。
\end{definition}

空带初始段经直接模拟：第 8 步首标记 $\mathbf 1^3$。

\begin{definition}[体制族与时代映射 $F$]
标记态按（形状，头位类）分属族 $\mathcal F\in\{A,\mathrm{g5B},\mathrm{g6o1},\mathrm{g6o2},\mathrm{g3},\mathrm{g1g1},\mathrm{g12},\mathrm{gE},\mathrm{g2L},\mathrm{g4A},\mathrm{g4},\mathrm{out},\mathrm{outR}\}$；每族的时代输出由下表给出（$T$ 逐字携带，除声明者）。
\end{definition}

\begin{center}\small
\begin{tabular}{llp{8.8cm}}
\toprule
族 & 形状（头位） & 时代输出\\
\midrule
$A$ & $1^L0^11^M0^51^RT$（$0^5$gap[1]） & $M\ge5$: $(L{+}3,M{-}4,R{+}1,T)$；$M{=}2\to$g5B$(L{-}2,2,R{+}1,T)$；$M\in\{1,4\}\to$g6$(L{+}3,R{+}1,T)$；$M{=}3\to$ 按尾首 $\to$ g1g1/g12/gE$(L{+}8,R,T)$\\
g5B & $1^L0^51^M0^21^RT$（gap[1]） & 见定理~\ref{thm:D1}\\
g6o1 & $1^L0^61^RT$（gap[1]） & 见定理~\ref{thm:dispatch}\\
g6o2 & $1^L0^61^RT$（gap[2]） & 见定理~\ref{thm:dispatch}\\
g3 & $1^L0^11^M0^31^RT$（$0^3$gap[1]） & 尾首 $(g2,S)\to A(L{+}3,M{+}R{-}4,\mathrm{g5},S{+}1,T)$；$(g1,S{\ge}2)\to A(L{+}3,M{+}R,\mathrm{g3},S{-}2,T)$；$(g1,1)$ 或尾空 $R\le2$ 变体 $\to$ \textbf{致死（定理~\ref{thm:lethal}）}；尾空 $R\ge1$ 按符号分派（定理~\ref{thm:dispatch}）\\
g1g1 & $1^L0^11^M0^11^RT$（$M$ 左缘） & $R\ge2\to$g4A；$R{=}1\to$\textbf{致死}\\
g12 & $1^L0^11^M0^21^RT$（$M$ 左缘） & $M{=}1\to$g6；$M\ge2\to$A/g5B\\
gE & $1^L0^11^M$（尾空） & $M{=}1\to$g6；$M{=}2\to$g5B(3,2)；$M{=}4\to$g6；$M\ge5\to$A\\
g2L & $1^L0^21^M0^21^RT$（$0^2$gap[1]） & $\to$g4；$M{=}1\to$\textbf{致死}\\
g4A & $1^L0^11^M0^41^RT$（gap[1]） & 按 $M{+}R{-}4$ 符号 $\to$ g3/g6；$R{=}1\to$\textbf{致死}\\
g4 & $1^L0^41^M0^21^RT$（gap[2]） & g4 自环 / A\\
out/outR & 头在块外 & outLeft$\to$g2L；outR $\to$ A\\
\bottomrule
\end{tabular}
\end{center}

\begin{theorem}[干链]\label{thm:dry}
设标记态 $1^L0^11^M0^51^RT\in A$、$M\ge5$、头 $h_0=L{+}M{+}2$。则该时代 = $\mathrm{Lap}^{\times M}+\mathrm{Fill}+\mathrm S+\mathrm{Dance}$，且输出标记态为
\[
1^{\,L+3}\,0^1\,1^{\,M-4}\,0^5\,1^{\,R+1}\,T .
\]
\end{theorem}
\begin{proof}
(1) \emph{$\mathrm{Lap}^{\times M}$}。归纳验证前件：第 $i$ 次 Lap 头在 $h_i=h_0-i+1$；$i=1$ 时 $\tau(h_0)=\tau(h_0{+}1)=0$（$0^5$ 内部）、$\tau(h_0{-}1)=0$（$0^5$ 首格）、$\tau(h_0{-}2)=1$（$1^M$ 末格）。若第 $i$ 次成立，其写入为 $\tau(h_i{-}2)\to0$、$\tau(h_i{+}1)\to1$；第 $i{+}1$ 次的头 $h_{i+1}=h_i-1$，前件格 $h_{i+1}{+}1=h_i$（$=0$，$C$ 所写）、$h_{i+1}{-}1=h_i{-}2$（$=0$，本步所写）、$h_{i+1}{-}2=h_i{-}3$（$=1$，$1^M$ 内部，$i\le M-1$ 时足长）。故 $M$ 次全为 Lap；侵蚀带 $[L{+}1,L{+}M]$ 全清 0，存入带 $[L{+}4,L{+}M{+}3]$ 逐格写 1，重叠区先蚀后存终值 $1$；结果 $1^L\,0^4\,1^M\,0^2$，头 $L{+}2$。
(2) \emph{Fill}。头 $L{+}2$，四前件格 $L{+}2,L{+}3,L{+}1,L$ 全为 $0$（$L$ 为原 $0^1$ 间隙格，未被侵蚀带触及），引理~\ref{lem:fill} 给 $1^{L+3}0^11^M0^2$、头 $L{+}4$。
(3) \emph{S}。头在 $1^M$ 左缘，$M$ 步至 $L{+}M{+}4$（$0^2$ 首格）。
(4) \emph{Dance}。头在 $0^2$ 的 gap[2] 变体，逐步字与引理~\ref{lem:dance} 相同（$F$ 扫描起点差 1）；净效应吃 $1^M$ 右缘 4 格、存入 $1^R$、新工作间隙 $0^5$，头回 gap[1]。得 $1^{L+3}0^11^{M-4}0^51^{R+1}T$。$T$ 全程未被触碰。
\end{proof}

\begin{theorem}[g5B 分派]\label{thm:D1}
设标记态 $1^L0^51^M0^21^RT\in\mathrm{g5B}$，$M\in\{2,3\}$，头在 $0^5$ gap[1]。则时代 = $\mathrm{Lap}^{\times L}+\mathrm{Fill}+\mathrm S+\mathrm{Dance}$，输出为：
$L\ge5$ 时 $A(3,L{-}4,M{+}1,R,T)$；$L=4$ 时 g6o2$(3,M{+}1)$；$L=3$ 时 g12 型 $(L{+}5,g_1,2,g_2,\dots)$；$L=2$ 时 \textbf{停机（H2）}；$L=1$ 时 outLeft。
\end{theorem}
\begin{proof}
$\mathrm{Lap}^{\times L}$ 与定理~\ref{thm:dry}(1) 同机制，侵蚀对象为左块 $1^L$（每 Lap $\tau(\cdot)\colon 1\to0$ 于 $h{-}2$ 即左块右缘、存入 $0^5$）；$L$ 次后左块尽蚀、$0^5$ 的右 $L$ 格为存入带：得 $0^3\,1^L\,0^{5-L}\,1^M0^21^RT$（$L\le5$）或其长存入变体，头至左块原左缘。Fill 重建左三格为 $1^3$：$1^3 0^1 1^L 0^2 1^M 0^2 1^RT$。S 越搬运块 $1^L$。Dance 自其右缘吃 4 格：余量 $L-4$ 决定分支——$L\ge5$ 余块 $1^{L-4}$ 成为新 $M$、存款并入 $1^M$ 得 $M{+}1$（输出 $A(3,L{-}4,M{+}1,R,T)$）；$L=4$ 恰尽、落点入 $0^6$ 型（g6o2$(3,M{+}1)$）；$L=3$ 舞链深入搬运块左侧间隙（g12 型）；$L=2$ 时 $F$ 扫描在冲刺块与存款块同耗尽后左入空白——$[y{-}1]=[y{-}2]=0$，由定理~\ref{thm:H}(ii) 停机；$L=1$ 同理左出（outLeft）。〔分派校验器对 $L\in\{1..12\}\times M\in\{2,3\}\times$ 边界 $R$ 全实例验证。〕
\end{proof}

\begin{theorem}[其余分派]\label{thm:dispatch}
\textbf{D2}（g6o1）：$L=1\to$out；$L=2\to$\textbf{停机（H2）}；$L=3\to$g2L$(8,2,1)$；$L=4\to$g6o2$(3,1)$；$L\ge5\to A(3,L{-}4,1,\text{前插}(g_1,R),T)$。
\textbf{D3}（g6o2，$L\ge3$）：$\to A(3,\,L{+}R{-}1,\ \text{尾块}{+}1,\ T)$。
\textbf{D4}（g3 终端，尾空）：$M'=M{+}R{-}4$：$M'\ge1\to A(L{+}3,M',\mathrm{g5},1,\varnothing)$；$M'=0\to$g6o2$(L{+}3,S{+}1)$；$M'=-1\to$gE$(M{+}R{+}4)$；$(M,R)=(1,1)\to$g5B$(1,2,1,\varnothing)$。
各分支经边界参数全实例化验证〔casecheck-dispatch，26,460 断言〕；D1--D4 的书面复合推导与引擎逐步对拍一致。
\end{theorem}

\section{停机刻画}

\begin{theorem}\label{thm:H}
$M$ 停机当且仅当下列二者之一发生：
\begin{enumerate}
\item[\textbf{(H1)}] 某时刻状态 $A$ 读格 $h$ 之 $0$，且 $\tau(h{+}1)=1$、$\tau(h{+}2)=0$；
\item[\textbf{(H2)}] 某时刻状态 $F$ 读格 $y$ 之 $1$，且 $\tau(y{-}1)=0$、$\tau(y{-}2)=0$。
\end{enumerate}
\end{theorem}
\begin{proof}
停机 ⟺ 到达 $C$ 读 $0$。转移表中以 $C$ 为目标者仅 $B0$（\texttt{1LC}）与 $F1$（\texttt{1LB}）。
(i) 经 $B0$：$C$ 所读格为 $B$ 刚写 $1$ 之格，值恒 $1$，不触发。故触发需 $B$ 读 $1$（\texttt{0RC}，擦除右行）后 $C$ 读 $[h{+}2]$。$B$ 读 $1$ 的来源：$A0$（\texttt{1RB} 写 1 右行）后 $B$ 读 $[h{+}1]$——触发当且仅当 $\tau(h{+}1)=1$（原值或 $A$ 所写）且 $\tau(h{+}2)=0$，即 H1。$F1$ 后 $B$ 读 $[y{-}1]$：若为 $0$ 则 $B0$ 写 1 左行、$C$ 读 $[y{-}2]$——触发当且仅当 $\tau(y{-}1)=0$ 且 $\tau(y{-}2)=0$，即 H2。$B$ 读 1 无其他入边（$D,E,F,A$ 的右行转移分别入 $E,F,A,B$ 之外的目标）。逆方向逐条构造即得停机路径。
\end{proof}

\section{扫掠守恒与致死集}

\begin{definition}[空尾扫掠]
尾空（$T=\varnothing$）标记态称空尾标记。$A$-族空尾标记的 $(L,M,R)$ 沿时代按定理~\ref{thm:dry} 演化，称一条\textbf{扫掠}。
\end{definition}

\begin{theorem}[扫掠守恒]\label{thm:sweep}
沿任一扫掠：$\mathrm{pop}=L{+}M{+}R$ 守恒；$Q:=L-3R$ 守恒；终端 $M\equiv M_0\pmod 4$（$M_0$ 为链首 $M$）。重置映射 $\Psi$：
\[
\Psi_{M4}\colon(L,M,R)\mapsto(3,\,\mathrm{pop}-1,\,1),\qquad
\Psi_{M3}\colon\mapsto(L{+}8,\,R{-}4,\,1),\qquad
\Psi_{M2}\colon R_0'=R+2,\quad \Psi_{M1}\colon R_0'=1 .
\]
\end{theorem}
\begin{proof}
守恒：$(L{+}3)+(M{-}4)+(R{+}1)=L{+}M{+}R$；$(L{+}3)-3(R{+}1)=L-3R$；$M\equiv M_0 \bmod 4$ 因每时代 $M\mapsto M-4$。$\Psi$ 各式由定理~\ref{thm:D1}/\ref{thm:dispatch} 对应终端支给出（$M4$ 支即 $M\in\{1,4\}\to$g6 族的参数展开；实例验证 $164/164$、$162/166$、$171/171$）。
\end{proof}

\begin{definition}[致死集 $P$]
$P$ = 满足下列之一的标记态参数格：(i) g1g1 族 $R=1$；(ii) g3 族尾首 $(g_1,1)$；(iii) g5B 族 $L=2$；(iv) g6o1 族 $L=2$；(v) g2L 族 $M=1$；(vi) g6o2 尾首 $(g_1,1)$；(vii) g4A 族 $R=1$（经 g1g1-R2 前插通路）。其闭包（对定理~\ref{thm:D1}/\ref{thm:dispatch} 前像取闭）共 \textbf{193 格}。
\end{definition}

\begin{theorem}[致死性]\label{thm:lethal}
每一格 $p\in P$ 被裸机验证为真停机：从对应构型出发，机器在 2--65 个时代内停机。特别地合成反例 $(L,M,R)=(7,2,1)$、$(8,2,2)$、$(21,1,1)$ 分别于 33、28、24 时代停机。反之，体制表（\S3）内其余一切格在 60--400 时代窗口内不停机〔census $1{,}152$ 格〕。
\end{theorem}

\begin{theorem}[归约]\label{thm:reduction}
以下等价：
\begin{enumerate}
\item $M$ 从空带启动停机；
\item 空带初始标记经定理~\ref{thm:dry}/\ref{thm:D1}/\ref{thm:dispatch} 的时代轨道 $(\mathfrak M_k)_{k\ge0}$ 满足 $\exists k:\ \mathfrak M_k\in P$。
\end{enumerate}
且经机器验证：轨道 $(\mathfrak M_k)$ 对 $k\le 2{,}830{,}885$（$3\times10^{10}$ 步）有 $\mathfrak M_k\notin P$。
\end{theorem}
\begin{proof}
(2)$\Rightarrow$(1)：定理~\ref{thm:lethal}。(1)$\Rightarrow$(2)：设停机，由定理~\ref{thm:H} 停机前最后状态属 H1 或 H2 之一。H1/H2 发生于某时代内部；对该时代所属族逐一审计落点（A/g5B/g6/g3/g2L/g4A/g4 族的冲刺块右间隙宽 $\ge2$，g1g1 族间隙 $=1$ 时 H1 $\iff R=1$，g3/g6o2 族 H1 $\iff$ 尾首 $(g_1,1)$，H2 的空白扫描几何仅 g5B-L2/g6o1-L2），H1/H2 的触发当且仅当该时代标记态 $\in P$〔落点审计 + 68 停机格普查〕。
\end{proof}

\section{入口封闭：一条已证之路}

\begin{theorem}[M4 入口封闭]\label{thm:m4}
若某空尾扫掠由 $M4$-重启（$\Psi_{M4}$）生成，则其后代 M3-终点的 $R$ 值不落致死谱系 $\{8\cdot4^n-1:n\ge0\}=\{1,7,31,127,\dots\}$。
\end{theorem}
\begin{proof}
$\Psi_{M4}$ 生成的扫掠上 $L=3R$ 精确成立（定理~\ref{thm:sweep}：$Q=0$ 初值），故 $\mathrm{pop}=4R$。M3-终点继承 $\mathrm{pop}\equiv0\pmod4$；而谱系值 $8\cdot4^n-1\equiv 3\pmod 4$。矛盾。
\end{proof}

初始标记、g6 入口与 M3$\to$M3 谱系自封闭性同理逐支核对〔dispatch 入口分类〕。未封闭的入口：M2-重启与 M1-重启。

\section{开放问题}\label{sec:open}

\begin{conjecture}[轨迹避让]\label{conj:avoid}
空带初始标记的时代轨道满足 $\forall k\ge0:\ \mathfrak M_k\notin P$。
\end{conjecture}

已验证：$k\le2{,}830{,}885$（$3\times10^{10}$ 步）；轨道相容格与 $P$ 的最小参数距离为 $3$（一例于 $t=34{,}922$：$(M1,L{=}48,R{=}20)$ 对 $(M1,L{=}51,R{=}20)$），且扫掠的对角运动（$\Delta L=3,\Delta R=1$ 同步）使异 $Q$ 格不可一步到达（定理~\ref{thm:sweep}）。已知否证的进攻路线：

\begin{enumerate}
\item \textbf{路线常数排除为假}：定理~\ref{thm:lethal} 的反例表明 M2/M1-重启路线常数并非全称避开谱系前置集。
\item \textbf{剩余类分离不存在}：对模 $3\times3$，可达剩余类与 $P$ 的 193 格全碰撞；且 M2/M1-重启的 $M_0'$ 不剩余类确定〔L6 型否定〕。
\item \textbf{线性不变量不充分}：$Q$ 守恒不蕴含避让——构型 g6o1$(11,1,\varnothing)$ 会从空尾制造尾首 $(g_1,1)$〔引擎复现〕；实轨 $L\equiv0\pmod3$ 使 $L\pm1$ 格不可达（143/193 格受此护城河保护）但此不变量为轨道观测，无书面归纳。
\end{enumerate}

\begin{remark}
范式定位：$(\mathfrak M_k)$ 是分段仿射系统的轨道，避让问题属分段仿射迭代的可达性类（Conway 1972 不可判定类的具体实例谱系；Kurtz--Simon 2007 独立性结果的适用范围见 paradigm-audit.md）。本机器停机命题的 ZFC 独立性不可排除。
\end{remark}

\section*{复跑}
\texttt{ocm.exe}（裸机 oracle）；\texttt{xval-macro.mjs}（忠实性三重对拍）；\texttt{census2/3.mjs}（体制普查）；\texttt{tailfront.mjs}（尾首演化）；\texttt{dispatch.mjs}（分派表）；\texttt{hclosure2.mjs}（族图不动点 294 节点）；\texttt{shape2.mjs 1e10}（深扫）；\texttt{lethal-db.json}（$P$ 库）。全部一行执行，见 \texttt{dossier.md}。

\end{document}
